%\chapter{\pysv\ 1.25}
% \vfill\pagebreak
\setlength{\footmarkwidth}{1.8em}
\pagewiselinenumbers
\dn 
\mll{YBH128INTRO}%
\llbl{p28_1}%
\mula{एवमवगतवाच्यवाचकसम्बन्धस्य \barfu योगिनः}%
\donote{YBH128INTRO}{{\dn एवमवगतवाच्यवाचकसम्बन्धस्य योगिनः}} \barfu %
\vrtsm{परमेश्वर}{\Tmpc\-\Td\-L\-M\-A\-\Me}{%
    \rdg{परमबर॰}{\Tmac}}%
प्रसादनं कथं क्रियत इति %\vrtsm{तत्प्र\-साद\-न\-साधन\-प्रतिपादना\-र्त्थ}{conj.}{\rdg{तद\-र्त्थ॰}{\all}}
तदर्त्थमाह---

\llbl{p28_2}%
\mll{YS128}\str{तज्जपस्तदर्त्थभावनम्॥२८॥}\donote{YS128}{{\dnb तज्जपस्तदर्त्थभावनम्॥२८॥}}

\llbl{p28_3}%
\mll{YBH128STRT}%
\mula{तस्ये}श्वर\mms वाचकस्य %
\mula{प्रणवस्या}र्द्धचतुर्म्मात्रस्य त्रि\mme मात्रस्य \barfu वा %
\mula{जपो} मनसो%
\vrtmm{पांशु वा}{\Me(em.)\vv{॰पांशु वा अा॰}{\Me}}{%
	\rdg{॰पांशूपा॰}{TLM}, 
	\rdg{॰पांशू॰}{A}}%
वर्त्त%
\vrtmm{न\-न्त\-ज्ज\-प}{conj.}{%
    \rdg{॰न\-ञ्ज\-प॰}{\Tmpc \Td\-L\-M\-A\-\Me{}\vv{॰नं ज॰}{\Td M\-A\-\Me}}, 
	\rdg{॰ञ्जपा॰}{\Tmac}}ः%
। %
\mula{तदर्त्थस्य} चेश्वरस्य प्रणवेन %
\vrt{वाचकेन}{T\Lpc M\-A\-\Me}{\rdg{वा\ccs क\cce चकेन}{L}} %
समर्प्पितस्य \barfu बुद्धौ \vrt{समा\-रो\-पि\-त\-स्य}{\Tmpc \Td L\-M\-A\-\Me}{\rdg{सम\ins ा\ine \ccs र्प्पि\-त\-स्य\cce रो\ccs त्\cce पि\-त\-स्य}{\Tm}} \mula{भावन}मभि\-ध्यान\-न्त\-दर्त्थ\-भावनं। कर्त्तव्य\mms मिति व\mme ाक्यशेषः। \mula{तदे}वमुभयं \mula{कुर्व्वतो} योगिन\mula{श्चित्तमेकाग्रं \vrt{स\-म्प\-द्य\-ते}{\Tmpc \Td L\-M\-A\-\Me}{\rdg{\ins सम्\ine पद्यते}{\Tm}}}॥\donote{YBH128STRT}{{\dn तस्य प्रणवस्य जपस्तदर्त्थस्य भावनं। तत्कुर्व्वतश्चित्तमेकाग्रं सम्पद्यते।}}

\llbl{p28_4}%
एकाग्र\-%
\vrtms{सम्पत्तेश्च}{em.}{%
	\rdg{॰सम्पत्तिश्च}{TLMA}, 
	\rdg{॰सम्पत्तिं च}{\Me}%
} %
तदाराधन%
\vrtmm{फल\-त्व}{em.}{%
	\rdg{॰फ\-ल॰}{\all}%
}%
न्दर्शयति %
\mll{PYS128QT}%
\mula{तथा चोक्तम् `%
\vrtsm{स्वा\-ध्या\-या}{\Tmpc\Td\-L\-M\-A\-\Me}{%
	\rdg{सा\-ध्या\-या॰}{\Tmac}%
}%
द्योगमा}%
\vrtms{\mula{तिष्ठेत्}}{TLMA}{%
	\rdg{॰(तिष्ठेत्)सीत}{\Me}%
} %
स्वा\-ध्यायात् प्रणव\-जपादी\-श्वरम् %
\vrtsm{प्रत्यवनत}{\Tmpc\-\Td\-L\-M\-A\-\Me}{%
	\rdg{प्रत्यवन\ins त\ine ॰}{\Tm}%
}%
चि\mms त्त\mme ः सन् योगमासीत। \barfu तदर्त्थमी%
\vrtms{श्व\-र\-न्ध्या\-येत्}{T\Lpc \Me{}\vv{॰रं ध्या॰}{\Td \Me}}{%
	\rdg{॰श्व\-र\-न्ध्या\ccs य\-\cce\-ये\-त्}{L}, 
	\rdg{॰श्व\-र\-स्य}{MA}%
}%
\fubah ॥%
\donote{PYS128QT}{{\dn तथा चोक्तम्---`स्वध्याया\-द्योगमा\-तिष्ठेत्}}

\llbl{p28_5}
\mll{PYS128QT2}
\vrtsm{तद}{TL}{\rdg{यत्तद॰}{MA}, \rdg{(रस्ययत्)तद॰}{\Me}}%
\vrtmm{र्त्थ\mula{ध्या}}{TM\-A\-\Me{}\vv{॰र्थ॰}{\Td\-M\-A\-\Me}}{%
  \rdg{॰र्त्थ\-न्ध्या॰}{L}%
}\mula{ना}\barfu%
\vrtmm{च्चा\-प्र\-च\-लि\-त}{L\-M\-A\-\Me}{%
  \rdg{॰च्च प्रचलित॰}{T}%
}%
\barfu मनाः %
\vrt{\mula{स्वाध्यायम्}}{T\Me{}\vv{॰यं}{\Td \Me}}{%
	\rdg{स्वा\-ध्या\-यः}{L\-M\-A}%
} %
प्रणव%
\vrtmm{म\barfu\mula{ ाम}}{TLM\Me}{%
	\rdg{॰मम॰}{A}%
}%
\mula{नेत्}%
\donote{PYS128QT2}{{\dn ध्या\-ना\-त्स्वा\-ध्या\-यमा\-मनेत्}} %
मनसाभिजपेत्। %
\vrt{मानसो}{TLM\Me}{%
	\rdg{मनसो}{A}%
} %
ऽभि\-जपः प्रशस्यते, \barfu ध्यानस्या%
\vrtmm{स\-न्न\-त\-र}{\Tmpc\Td\-L\-M\-A\-\Me}{%
	\rdg{॰त्तन्नतर॰}{\Tmac}%
}%
त्वात्, मा विषय%
\vrtmm{प्रवण}{T\Me}{%
	\rdg{॰प्र\-ण\-व॰}{LMA}%
}%
चित्तो %
\vrtsm{भू}{\Tm\-\Tdpc\-L\-\Mpc\-A\-\Me}{%
	\rdg{\ccs ऽ\cce भू॰}{\Td\-M}%
}%
दिति॥%

\begin{messwithp}{1000}{10pt}

\llbl{p28_6}
\mll{PYS128QT3}%
एवं \mula{स्वाध्याययोगसम्पत्या}%
\donote{PYS128QT3}{{\dn स्वाध्याय\-योग\-स\-म्प\-त्या}}%
---तद्विरोधिना प्रत्ययान्तरेणान%
\vrtms{भि\-घात}{\all}{%
	\rdg{भिघातः}{\Tm L (note the presence of \emph{visarga})}%
} %
\barfu स्वाध्याययोग\-सम्प\-त्तिः॥%

\end{messwithp}

\llbl{p28_7}
\mll{PYS128QT4}%
तथा च \vrtsm{प्रणव}{\Tmpc \Td L\-M\-A\-\Me}{\rdg{प्र\ccs प्र\cce णव॰}{\Tm}}जपपरमेश्वरध्यानसम्पत्या \mula{पर \barfu आत्मा} %
% \donote{PYS128QT}{%
% 	{\dn तथा चोक्तम्---स्वध्याया\-द्योगमा\-तिष्ठेद्ध्या\-ना\-त्स्वा\-ध्या\-यमा\-मनेत्। स्वाध्याय\-योग\-स\-म्प\-त्या पर \barfu आत्मा प्रकाशते॥ इति}%
% } %
परमेष्ठी \mula{प्रकाशते'} %
योगिन %
\vrt{\mula{इति}॥२८॥}{\Me}{\rdg{इति}{TLMA\vv{इति।}{\Td MA} (no end of section marker)}}%
\donote{PYS128QT4}{{\dn पर \barfu आत्मा प्रकाशते' इति}} %

\nolinenumbers

\vfil
\reseteditionparcounter
